% GENERATED FILE. Edit canonical lessons, then run npm run generate:books.
% canonical-source-hash: fnv1a64:af1766a84f8bd4e8
% canonical-lessons: RU-C366-vybory, RU-C366-golosovat, RU-C366-partiya, RU-C366-protest, RU-C366-protestovat

\chapter{Elections}
\label{ch:ru-elections}
\hlchaptermodality{366}

\begin{chapteropening}
\textbf{By the end of this chapter:} \emph{I can talk about an election, say whom I vote for, speak about a party and a protest, and say that people protest.}

Complete the last lesson of chapter 366: I can talk about an election, say whom I vote for, speak about a party and a protest, and say that people protest.
\end{chapteropening}

\section[\texorpdfstring{\ru{выборы} (výbory)}{выборы (výbory)}]{\ru{выборы} (výbory) — an election}

\label{lesson:RU-C366-vybory}



\begin{quote}
Before the new one: say the Russian for politics, then the Russian for a politician.
\end{quote}

\subsection*{You'll want to know: \ru{выборы}}
\textbf{\ru{выборы}} (\emph{výbory}) — \textquotedblleft{}an election\textquotedblright{}.

\begin{quote}
\ru{Выборы} \ru{будут} \ru{в} \ru{сентябре}.
\end{quote}

\begin{quote}
\emph{Výbory búdut v sentyabré.}
\end{quote}

\begin{quote}
The election will be in September.
\end{quote}

Plural only. \emph{\ru{Выборы} \ru{президента}} — a presidential election.

\begin{cousinweb}
\textbf{\ru{выборы}} is built from \emph{\ru{вы}-} \textquotedblleft{}out\textquotedblright{} and the root of \emph{\ru{брать}} \textquotedblleft{}to take\textquotedblright{}, as in \emph{\ru{выбирать}} \textquotedblleft{}to choose\textquotedblright{}: an election picks people out.
\end{cousinweb}

\begin{grammarlens}[title={What you've built}]
One more word for telling, explaining and describing at B1.
\end{grammarlens}

\subsection*{Your turn}
\begin{itemize}
\raggedright
  \item \emph{Say it:} \emph{výbory}
  \item \emph{Say it:} \emph{výbory}, once more
  \item \emph{Say it:} a sentence of your own with \emph{výbory}
  \item \emph{Recall:} say the Russian for politics, then the Russian for a politician, then say \emph{výbory} again
\end{itemize}

\begin{tcolorbox}[breakable,colback=teal!4,colframe=teal!35!black,title={Before you move on}]
What does \ru{выборы} mean? (\textquotedblleft{}An election\textquotedblright{}.) Say it once more.
\end{tcolorbox}
\section[\texorpdfstring{\ru{голосовать} (golosovát)}{голосовать (golosovát)}]{\ru{голосовать} (golosovát) — to vote}

\label{lesson:RU-C366-golosovat}



\begin{quote}
Before the new one: say the Russian for a politician, then the Russian for an election.
\end{quote}

\subsection*{You'll want to know: \ru{голосовать}}
\textbf{\ru{голосовать}} (\emph{golosovát}) — \textquotedblleft{}to vote\textquotedblright{}.

\begin{quote}
\ru{Мы} \ru{голосовали} \ru{за} \ru{нового} \ru{мэра}.
\end{quote}

\begin{quote}
\emph{My golosováli za nóvogo méra.}
\end{quote}

\begin{quote}
We voted for a new mayor.
\end{quote}

Imperfective; the present is \emph{\ru{голосую}, \ru{голосуешь}}; the perfective is \emph{\ru{проголосовать}}. \emph{\ru{Голосовать} \ru{за}} + accusative, \emph{\ru{против}} + genitive.

\begin{cousinweb}
\textbf{\ru{голосовать}} is built on \emph{\ru{голос}} \textquotedblleft{}voice\textquotedblright{}, inherited through Old East Slavic \emph{golosŭ} from Proto-Slavic \emph{*golsŭ} \textquotedblleft{}voice\textquotedblright{} (EtymDB). The \emph{-\ru{оло}-} is the native Russian shape: to vote is to give your voice.
\end{cousinweb}

\begin{grammarlens}[title={What you've built}]
One more word for telling, explaining and describing at B1.
\end{grammarlens}

\subsection*{Your turn}
\begin{itemize}
\raggedright
  \item \emph{Say it:} \emph{golosovát}
  \item \emph{Say it:} \emph{golosovát}, once more
  \item \emph{Say it:} a sentence of your own with \emph{golosovát}
  \item \emph{Recall:} say the Russian for a politician, then the Russian for an election, then say \emph{golosovát} again
\end{itemize}

\begin{tcolorbox}[breakable,colback=teal!4,colframe=teal!35!black,title={Before you move on}]
What does \ru{голосовать} mean? (\textquotedblleft{}To vote\textquotedblright{}.) Say it once more.
\end{tcolorbox}
\section[\texorpdfstring{\ru{партия} (pártiya)}{партия (pártiya)}]{\ru{партия} (pártiya) — a party (in politics); a game}

\label{lesson:RU-C366-partiya}



\begin{quote}
Before the new one: say the Russian for an election, then the Russian for to vote.
\end{quote}

\subsection*{You'll want to know: \ru{партия}}
\textbf{\ru{партия}} (\emph{pártiya}) — \textquotedblleft{}a party (in politics); a game\textquotedblright{}.

\begin{quote}
\ru{Он} \ru{член} \ru{этой} \ru{партии}.
\end{quote}

\begin{quote}
\emph{On chlen étoy pártii.}
\end{quote}

\begin{quote}
He is a member of this party.
\end{quote}

Feminine. Also one game of chess: \emph{\ru{партия} \ru{в} \ru{шахматы}}.

\begin{cousinweb}
\textbf{\ru{партия}} was borrowed from Polish \emph{partia}, from French \emph{partie} \textquotedblleft{}a part\textquotedblright{}, from \emph{partir} \textquotedblleft{}to divide\textquotedblright{}, from Latin \emph{partīre}, built on \emph{pars} \textquotedblleft{}a part, a piece\textquotedblright{} (EtymDB).
\end{cousinweb}

\begin{grammarlens}[title={What you've built}]
One more word for telling, explaining and describing at B1.
\end{grammarlens}

\subsection*{Your turn}
\begin{itemize}
\raggedright
  \item \emph{Say it:} \emph{pártiya}
  \item \emph{Say it:} \emph{pártiya}, once more
  \item \emph{Say it:} a sentence of your own with \emph{pártiya}
  \item \emph{Recall:} say the Russian for an election, then the Russian for to vote, then say \emph{pártiya} again
\end{itemize}

\begin{tcolorbox}[breakable,colback=teal!4,colframe=teal!35!black,title={Before you move on}]
What does \ru{партия} mean? (\textquotedblleft{}A party (in politics); a game\textquotedblright{}.) Say it once more.
\end{tcolorbox}
\section[\texorpdfstring{\ru{протест} (protést)}{протест (protést)}]{\ru{протест} (protést) — a protest}

\label{lesson:RU-C366-protest}



\begin{quote}
Before the new one: say the Russian for to vote, then the Russian for a party.
\end{quote}

\subsection*{You'll want to know: \ru{протест}}
\textbf{\ru{протест}} (\emph{protést}) — \textquotedblleft{}a protest\textquotedblright{}.

\begin{quote}
\ru{Студенты} \ru{вышли} \ru{на} \ru{протест}.
\end{quote}

\begin{quote}
\emph{Studénty výshli na protést.}
\end{quote}

\begin{quote}
The students came out to protest.
\end{quote}

Masculine. \emph{\ru{В} \ru{знак} \ru{протеста}} — as a sign of protest.

\begin{cousinweb}
\textbf{\ru{протест}} is usually traced, through western languages, to Latin \emph{prōtestārī} \textquotedblleft{}to declare in public\textquotedblright{}, built on \emph{testis} \textquotedblleft{}a witness\textquotedblright{} (Lewis \& Short).
\end{cousinweb}

\begin{grammarlens}[title={What you've built}]
One more word for telling, explaining and describing at B1.
\end{grammarlens}

\subsection*{Your turn}
\begin{itemize}
\raggedright
  \item \emph{Say it:} \emph{protést}
  \item \emph{Say it:} \emph{protést}, once more
  \item \emph{Say it:} a sentence of your own with \emph{protést}
  \item \emph{Recall:} say the Russian for to vote, then the Russian for a party, then say \emph{protést} again
\end{itemize}

\begin{tcolorbox}[breakable,colback=teal!4,colframe=teal!35!black,title={Before you move on}]
What does \ru{протест} mean? (\textquotedblleft{}A protest\textquotedblright{}.) Say it once more.
\end{tcolorbox}
\section[\texorpdfstring{\ru{протестовать} (protestovát)}{протестовать (protestovát)}]{\ru{протестовать} (protestovát) — to protest}

\label{lesson:RU-C366-protestovat}



\begin{quote}
Before the new one: say the Russian for a party, then the Russian for a protest.
\end{quote}

\subsection*{You'll want to know: \ru{протестовать}}
\textbf{\ru{протестовать}} (\emph{protestovát}) — \textquotedblleft{}to protest\textquotedblright{}.

\begin{quote}
\ru{Жители} \ru{протестуют} \ru{против} \ru{стройки}.
\end{quote}

\begin{quote}
\emph{Zhíteli protestúyut prótiv stróyki.}
\end{quote}

\begin{quote}
The residents are protesting against the building work.
\end{quote}

Imperfective; the present is \emph{\ru{протестую}, \ru{протестуешь}}. \emph{\ru{Протестовать} \ru{против}} + genitive.

\begin{cousinweb}
\textbf{\ru{протестовать}} is built on \emph{\ru{протест}} with the verb suffix \emph{-\ru{овать}}; the noun is usually traced to Latin \emph{prōtestārī} \textquotedblleft{}to declare in public\textquotedblright{}, on \emph{testis} \textquotedblleft{}a witness\textquotedblright{} (Lewis \& Short).
\end{cousinweb}

\begin{grammarlens}[title={What you've built}]
One more word for telling, explaining and describing at B1.
\end{grammarlens}

\subsection*{Your turn}
\begin{itemize}
\raggedright
  \item \emph{Say it:} \emph{protestovát}
  \item \emph{Say it:} \emph{protestovát}, once more
  \item \emph{Say it:} a sentence of your own with \emph{protestovát}
  \item \emph{Recall:} say the Russian for a party, then the Russian for a protest, then say \emph{protestovát} again
\end{itemize}

\begin{tcolorbox}[breakable,colback=teal!4,colframe=teal!35!black,title={Before you move on}]
What does \ru{протестовать} mean? (\textquotedblleft{}To protest\textquotedblright{}.) Say it once more.
\end{tcolorbox}
